%% Section 5: Conclusion
\section{Conclusion}
\label{sec:conclusion}

This paper set out to formalize a pedagogical scenario for mobile AR-enhanced vocabulary learning using the Unified Modeling Language. The motivation was twofold: the persistent gap between the pedagogical potential of AR in education and the absence of rigorous, implementation-ready models that integrate AR with intelligent tutoring, OCR-based content acquisition, and gamification; and the demonstrated effectiveness of UML-driven approaches in bridging pedagogical design intent and technical specification~\cite{Booch2005,Ouariach2023,Ouariach2026}.

Through four complementary UML diagrams---class, use case, activity, and sequence---we translated the abstract educational objectives of the AR-DeC (Augmented Reality Didactic Companion) system into a concrete, traceable architectural model. The class diagram specifies fifteen classes organized into five functional groups (Core Actors, Learning Content, Student Modeling, AI \& AR, and Gamification), each encapsulating a coherent set of pedagogical responsibilities. The use case diagram delineates the interactions of three actors---Student, Teacher, and System---through primary use cases, \texttt{<<include>>} relationships for mandatory sub-behaviors, and \texttt{<<extend>>} relationships for conditional help-seeking pathways. The activity diagram captures the dynamic workflow of a complete learning session, from application launch and streak verification through OCR scanning, adaptive explanation generation via Bloom's taxonomy~\cite{Bloom1956,Krathwohl2002}, AR overlay rendering, gamified quizzing, and backend synchronization. The sequence diagram specifies the chronological message exchange between six key lifelines during a representative scan-learn-quiz cycle, using \texttt{par}, \texttt{opt}, and \texttt{alt} interaction fragments to model concurrency, optionality, and conditional branching.

The resulting framework makes three principal contributions. First, it demonstrates that UML can serve not merely as a documentation tool but as an instructional design instrument, encoding pedagogical decisions---such as the mapping from BKT mastery estimates to Bloom's taxonomy levels~\cite{Corbett1994}, or the tiered escalation from conversational agent to intelligent tutor to human teacher---in a formal, verifiable notation. Second, the fifteen-class architecture provides a clear separation of concerns that facilitates independent development, testing, and future extension of each component. Third, the dual-perspective approach---combining static structural views with dynamic behavioral views---ensures that both the \emph{what} and the \emph{how} of the pedagogical scenario are fully specified, reducing the risk of pedagogical drift during implementation~\cite{Rumbaugh2004}.

Although the framework remains at the conceptual stage and requires empirical validation, it provides a robust foundation for the next phases of development. The UML diagrams serve as a precise specification from which a functional prototype can be systematically constructed, and the explicit traceability between educational requirements and architectural components ensures that the implemented system will faithfully reflect the pedagogical intent~\cite{Ouariach2025,Ouariach2026}.

Looking ahead, our immediate next step is the implementation of a working AR-DeC prototype using the technology stack specified in the architecture (React Native with Expo, on-device OCR via ML Kit, and AR rendering), followed by a rigorous empirical evaluation through Design-Based Research methodology~\cite{Crompton2017}. Such evaluation will assess the impact of the integrated system on vocabulary acquisition outcomes, learner engagement, and long-term knowledge retention, comparing the full AR-DeC experience against reduced configurations to isolate the contribution of each component. Beyond AR-DeC itself, the UML-driven modeling approach presented here is generalizable: it can be adapted to other educational domains---science laboratories, historical site exploration, medical training---wherever the combination of real-world scanning, intelligent adaptation, immersive visualization, and motivational scaffolding can enhance learning. We hope that this work encourages researchers and developers to adopt formal modeling as a foundational step in the design of pedagogically grounded AR educational systems.
